% Generated from locale/tr/content/turing-machines/undecidability/introduction.tex; do not hand-edit.


\olfileid[tr]{tur}{und}{int} 
\olsection{Giriş}

Her fonksiyonun, hatta her aritmetik fonksiyonun hesaplanabilir olamayacağı
açık görünebilir. Sayıları fazladır ve davranışları fazlasıyla karmaşıktır.
Örneğin radyoaktif parçacıkların bozunmasından ya da başka kaotik veya
rastlantısal davranışlardan tanımlanan fonksiyonlar böyledir. Belirli bir
andan başlayarak bir saniyelik aralıkları saydığımızı ve $f(n)$ fonksiyonunu,
bu ilk andan sonraki $n$. bir saniyelik aralıkta evrende bozunan parçacık
sayısı olarak tanımladığımızı varsayalım. Bu, hesaplamayı hiçbir zaman
umamayacağımız bir fonksiyon adayı gibi görünür.

Ne var ki böyle fonksiyonların nasıl hesaplanacağını tasarlayamamak başka,
gerçekten hesaplanamaz olduklarını kanıtlamak başkadır. Umutsuzca karmaşık
görünen fonksiyonlar bile soyut bir anlamda hesaplanabilir olabilir. Örneğin
evrenin zamanda sonlu olduğunu, bazı kozmoloji kuramlarının öngördüğü gibi çok
uzak bir gelecekte bir gün tek bir noktaya çökeceğini varsayalım. O zaman bu
ilk andan başlayarak $f(n)$'nin tanımlı olduğu yalnızca sonlu (ama inanılmaz
büyük) sayıda saniye vardır. Yalnızca sonlu sayıda girdide tanımlı her
fonksiyon hesaplanabilir: Çıktıları büyük bir tabloda listeleyebilir ya da çok
büyük bir Turing makinesi durum geçiş diyagramında kodlayabiliriz.

Değerleri evet/hayır sorularının yanıtlarını veren özel fonksiyonlarla sık
sık ilgileniriz. Örneğin ``$n$ bir asal sayı mıdır?'' sorusu
\[
\fn{isprime}(n) = \begin{cases}
  1 & \text{$n$ asal ise}\\
  0 & \text{aksi durumda}
  \end{cases}
\]
fonksiyonuyla ilişkilidir. İlişkili $1/0$ değerli fonksiyon etkili biçimde
hesaplanabiliyorsa bir evet/hayır sorusuna \emph{etkili biçimde karar
verilebildiğini} söyleriz.

Etkili biçimde hesaplanamayan fonksiyonların ya da etkili biçimde karar
verilemeyen problemlerin bulunduğunu matematiksel olarak kanıtlamak için
belirli bir hesaplama modelini sabitlemek ve bu modelin hesaplayamadığı
fonksiyonlar ya da karar veremediği problemler bulunduğunu göstermek gerekir.
Örneğin her fonksiyonun Turing makinelerince hesaplanamayacağını ve her
probleme Turing makinelerince karar verilemeyeceğini gösterebiliriz. Ardından,
yalnızca Turing makinelerinin her fonksiyonu hesaplayacak güçte olmadığı değil,
hiçbir etkili sürecin bunu yapamayacağı sonucuna varmak için Church--Turing
tezine başvurabiliriz.

Bu tür olumsuz sonuçları kanıtlamanın anahtarı, Turing makinelerinin
kendilerine sayı atayabilmemizdir. Bunu yapmanın en kolay yolu, belki Turing
makinelerini ve programlarını yazmanın belirli bir biçimini sabitleyip sonra
onları sistematik olarak listeleyerek, makineleri numaralandırmaktır. Bunun
yapılabildiğini gördüğümüzde Turing-hesaplanamaz fonksiyonların varlığı basit
kardinalite değerlendirmelerinden çıkar: $\Nat$'ten $\Nat$'e (hatta yalnızca
$\Nat$'ten $\{0,1\}$'e) fonksiyonlar kümesi !!{nonenumerable}, fakat bütün
Turing makinelerini numaralandırabildiğimiz için Turing-hesaplanabilir
fonksiyonlar kümesi yalnızca !!{denumerable}.

Sırasıyla hesaplanamaz ve karar verilemez olduklarını kanıtlayabileceğimiz
\emph{belirli} fonksiyonlar ve problemler de tanımlayabiliriz. Bunlardan biri
\emph{Durma Problemi} denen problemdir. Turing makineleri, yönergeleri
listelenerek sonlu biçimde betimlenebilir. Bir Turing makinesinin böyle bir
betimi, yani Turing makinesi programı, elbette başka bir Turing makinesine
girdi olarak verilebilir. Dolayısıyla başka Turing makineleri hakkındaki
sorulara karar veren Turing makinelerini ele alabiliriz. Özellikle ilginç bir
soru şudur: ``Verilen Turing makinesi $n$ girdisinde başlatıldığında sonunda
durur mu?'' Bu soruya karar verebilen bir Turing makinesi bulunsa iyi olurdu;
onu, Turing makinelerinin sonsuz döngülere vb. takılmamasını sağlayan bir
kalite denetim Turing makinesi olarak düşünün. Turing'in kanıtladığı ilginç
olgu, böyle bir Turing makinesinin bulunamayacağıdır. Bir Turing makinesi~$M$
ile bir $n$ sayısının betiminden oluşan girdide başlatıldığında, $M$ makinesi
$n$ girdisinde başlatılırsa durup durmayacağına göre her zaman $1$ ya da $0$
çıktısıyla duran tek bir Turing makinesi bulunamaz.

Belirli karar verilemez problem örneklerine sahip olduktan sonra bunları başka
problemlerin de karar verilemez olduğunu göstermek için kullanabiliriz.
Örneğin ünlü bir karar verilemez problem, ``Birinci dereceden
!!{formula}~$!A$ geçerli midir?'' sorusudur. Birinci dereceden
!!{formula}~$!A$ girdi olarak verildiğinde $!A$'nın geçerli olup olmamasına
göre $1$ ya da $0$ çıktısıyla duracağı güvence altında olan hiçbir Turing
makinesi yoktur. Tarihsel olarak bu problemi etkili biçimde çözecek bir süreç
bulma sorusuna yalnızca ``karar problemi'' denmiştir; bu nedenle karar
probleminin çözülemez olduğunu söyleriz. Turing ile Church bu sonucu yaklaşık
aynı zamanda birbirlerinden bağımsız olarak kanıtlamışlardır; bu yüzden sonuca
Church--Turing Teoremi de denir.

